once from that of the corresponding squares (7.10.0), taking account of the transitivity of inverse images (IV 2.12).

\sgaStableTarget{sga6.e10.rem.7-10-3}\textbf{Remark 7.10.3.} When @@M0@@ is proper, the commutativity of (7.10.1) immediately implies that of the following diagram, obtained from it by passage to the associated graded groups:
\sgaStableTarget{sga6.e10.eq.7-10-1bis}@@M1@@
Likewise, when @@M2@@ and @@M3@@ have regular fibers and are separated, the commutativity of (7.10.2) implies that of
\sgaStableTarget{sga6.e10.eq.7-10-2bis}@@M4@@
When @@M5@@ and @@M6@@ are moreover quasi-projective and smooth over @@M7@@, one also knows that @@M8@@ and @@M9@@ are compatible with the filtrations by codimension of supports, and one obtains a commutative diagram like (7.10.2 bis), but with the @@M10@@ replaced by the @@M11@@ (cf. Exposé 0, 2.1). One can also note that, when @@M12@@ and @@M13@@ are quasi-projective over @@M14@@ with regular fibers and separated, @@M15@@ and @@M16@@ are automatically complete-intersection morphisms (VIII 1.1), and one obtains a commutative diagram analogous to (7.10.1 bis), with the @@M17@@ replaced by the @@M18@@, taking account of the compatibility result (VIII 3.2).
